
\begin{center}
    \bfseries \zihao{3} 毕业论文（设计）原创性声明
\end{center}

{
    \zihao{4}
    \setlength{\parindent}{2em}

    本人郑重声明：所呈交的毕业论文（设计）是本人在导师的指导下，严格按照学校和学院有关规定完成的，不存在学术不端行为。除文中特别加以标注和致谢的地方外，本论文（设计）不包含其他个人或集体已经发表或撰写过的研究成果，也不包含为获得\bful{~浙江大学~}或其他教育机构的学位或证书而使用过的材料。对本研究做出贡献的个人和集体，均已在文中明确说明并表示谢意。

    \vspace{0em}

    \noindent
    \hspace{6em} 毕业论文（设计）作者签名：

    \vspace{0.5em}

    \noindent
    \hspace{14em} 签字日期：\hspace{4em} 年 \quad 月 \quad 日
}

\vspace{3em}

\begin{center}
    \bfseries \zihao{3} 毕业论文（设计）使用授权书
\end{center}

{
    \zihao{4}
    \setlength{\parindent}{2em}
    \sloppy                          % 放宽断行容忍度
    \setlength{\emergencystretch}{3em} % 实在断不了就拉伸这么多

    本人完全了解\bful{~浙江大学~}有权保留并向国家有关部门或机构送交本论文（设计）的复印件和电子版，
    许本论文（设计）被查阅和借阅。本人授权\bful{~浙江大学~}可以将本论文（设计）
    的全部或部分内容编入有关数据库进行检索和传播，可以采用影印、缩印或扫描等复制手段保存、汇编本论文（设计）。

    （涉密毕业论文（设计）在解密后适用本授权书。）

    \vspace{2em}

    \begin{center}
        \begin{tabularx}{\linewidth}{>{\fangsong}l>{\fangsong}X>{\fangsong}l>{\fangsong}X}
            毕业设计作者签名： & ~ & 导师签名： & ~ \\

            \vspace{0.4em}

            签字日期： & 年 \quad 月 \quad 日 & 签字日期： & 年 \quad 月 \quad 日 \\
        \end{tabularx}
    \end{center}
}
